\documentclass{article}
\usepackage{amsmath,hyperref}
\title{Exact phase recovery under substitutions for an optical codebook}
\begin{document}
\maketitle

\paragraph{Result.}
For the 84-word, 11-bit alphabet in the
\href{../Optical_Frame_Acquisition.tex}{parent study}, the exact
worst-case number of consecutive samples needed to identify word
phase with at most \(e\) substitutions is \(17+11e\), for every
integer \(e\geq 0\). This uses the entire observation, not just
the newest candidate frame. It assumes a single contiguous stream
of arbitrary words from this exact codebook, accurate bit cadence,
unknown initial word phase, and no bit insertions, deletions or
mid-observation timing changes. It retains the sender format and
word count. It does not give the receiver an epoch, a fresh receipt,
or a bound in physical time.

\paragraph{Why the bound is exact.}
Make each proper prefix of an admissible codeword a state,
including the empty prefix. An admissible bit extends it; the
last bit of a frame returns to the empty prefix. The depth is the
number of bits already consumed in a frame. Every initial state
is possible, and each retained state can continue. In the product
of two copies, start from all pairs of different depths, keeping
just one orientation at time zero. Each step costs one if the two
bits differ, zero otherwise. Let \(D_n\) be the vector of minimum
costs to reach each state pair after \(n\) bits, with infinity for
an unreachable pair. The least finite entry \(d_n\) is the minimum
Hamming distance between any two length-\(n\) clean trajectories
of different initial phases. The constant difference of depths
modulo 11 keeps those phases distinct.

Two such trajectories can explain the same observation with at
most \(e\) substitutions each exactly when their distance is at
most \(2e\). Necessity is the triangle inequality. Sufficiency
follows by splitting their differing positions between the two
explanations, choosing at each differing position one of the two
bits. Therefore phase is determined at \(n\) precisely when
\(d_n>2e\). The minimum cannot decrease with length because
each longer pair has an admissible prefix pair with the same
phase difference.

The finite replay checks \(D_{27}=D_{16}+2\) coordinate by
coordinate, including equality of the reachable support under
the convention \(\infty+2=\infty\). It checks 349 trie states,
48,896 initially possible pairs and 48,896 finite coordinates
in the two equal-support vectors. The next-step map takes
minima of incoming costs plus the cost of that edge. It
commutes with adding the same constant to all finite entries
and keeping infinity. Applying that map repeatedly proves
\(D_{n+11}=D_n+2\) for every \(n\geq16\). The checked minima
are \(d_{16}=0\) and \(d_{17}=1\), so
\[
  d_{16+11e}=2e,\qquad d_{17+11e}=2e+1.
\]
This proves both necessity and sufficiency for every budget,
without extrapolating a list of finite-budget outputs. A separate
fixed-phase checker splits the received run at the proposed
boundaries and takes minimum distances to matching partial or full
codewords. It verifies a 27-bit one-substitution witness with
exactly phases 2 and 5, counted as bits already consumed.
That checker is also a constructive decoder at the threshold.

\paragraph{Identity of a full word.}
At one substitution, 28 samples identify phase. To identify the
actual content of some wholly observed frame takes 32 samples in
the worst case. At 32, even the latest next boundary leaves two
complete frames. At least one is clean, and the other cannot
masquerade as a different codeword after one substitution because
distinct codewords have even weight and thus distance at least
two. At 31, a cut after only the first bit of a frame gives a
10-bit suffix, one complete frame, and a 10-bit prefix. Two
admissible words at distance two have a one-flip midpoint;
placing it in that complete frame makes the only full word
ambiguous even when the phase is supplied. The replay checks
the concrete binary witness.

At two or more substitutions, there is no finite uniform bound
for identifying the content of any particular complete position
in an arbitrary stream. Let the observation be a repetition
of one of the distance-two words. For every chosen complete
position, a clean trajectory with the alternative word just
at that position fits the same observation. It uses only two
substitutions. Although a sufficiently long observation
contains some clean frames, the receiver need not be able to
identify one as correct. Repeated emission of the same needed
receipt by contract, or a code with greater separation, is a
different problem and can change this result.

\paragraph{Use in the project.}
This is exact for the declared bit model and unchanged alphabet.
It strengthens the reference from a sliding-word check to
a whole-observation distance decoder; that direct decoder
agrees with the phase criterion. It supplies a bound on
samples, not on camera delivery, physical clearance,
message authentication, or freshness. A physical transfer
must independently test whether the source really meets the
error and timing model and whether a decoded payload can
justify the receiver's action.
\end{document}
